\footnotesize
\setlength{\tabcolsep}{4pt}
\begin{tabular}{llrrrrr}
\toprule
Tercile & Range & Coefficient & Bootstrap $p$ & Quarters & \multicolumn{2}{c}{Mean provisioning} \\
\cmidrule(lr){6-7}
 & & (clustered SE) & & coop.\ / banks & Coop. & Banks \\
\midrule
\multicolumn{7}{l}{\textit{Panel A. Terciles of the Basel ratio residualised on log assets and quarter fixed effects}} \\
Bottom & [-40.8, -9.7] & -0.0105 (0.0095) & 0.290 & 1,375 / 796 & 0.053 & 0.063 \\
Middle & [-9.7, -0.6] & 0.0007 (0.0104) & 0.936 & 1,181 / 1,105 & 0.054 & 0.052 \\
Top & [-0.6, 165.1] & 0.0109 (0.0139) & 0.446 & 389 / 1,605 & 0.063 & 0.047 \\
\midrule
\multicolumn{7}{l}{\textit{Panel B. Terciles of the unconditional Basel ratio}} \\
Bottom & [9.8, 16.8] & -0.0081 (0.0086) & 0.355 & 657 / 1,598 & 0.048 & 0.058 \\
Middle & [16.8, 23.9] & 0.0178*** (0.0056) & 0.001 & 1,344 / 881 & 0.054 & 0.039 \\
Top & [23.9, 189.4] & -0.0007 (0.0127) & 0.951 & 944 / 1,027 & 0.060 & 0.057 \\
\bottomrule
\end{tabular}
\begin{tablenotes}[flushleft]\footnotesize
\item Coefficient on the cooperative dummy in a regression of the provisioning ratio on the dummy, log assets and quarter fixed effects, estimated separately within each tercile of the analysis sample. Terciles are cut on the pooled sample of cooperatives and banks: in Panel A on the residual of the Basel ratio from a regression on log assets and quarter fixed effects, in Panel B on the ratio itself; the range column gives the tercile's cut points on that variable. Standard errors clustered by institution in parentheses; stars and the $p$-value from the restricted wild cluster bootstrap with 1,999 replications (*** $p<0.01$, ** $p<0.05$, * $p<0.10$). Quarters are those with a non-missing provisioning ratio.
\end{tablenotes}
